\documentclass[10pt, a4paper]{article}

\usepackage[a4paper, landscape, left=9mm, right=9mm, top=11mm, bottom=13mm]{geometry}
\usepackage{fontspec}
\usepackage[greek, bidi=basic, provide=*]{babel}
\babelprovide[import, onchar=ids fonts]{english}
\babelfont{rm}{Noto Sans}

\usepackage{enumitem}
\setlist[itemize]{label=-}

\usepackage[table]{xcolor}
\usepackage{longtable}
\usepackage{array}
\usepackage{ragged2e}
\usepackage{fancyhdr}
\usepackage{lastpage}
\usepackage[hidelinks]{hyperref}

% Χρωματική παλέτα
\definecolor{primaryblue}{HTML}{1B365D}
\definecolor{accentblue}{HTML}{2C5282}
\definecolor{rowalt}{HTML}{F4F7FA}
\definecolor{bordergray}{HTML}{CBD5E1}
\definecolor{critred}{HTML}{9B1C1C}
\definecolor{highamber}{HTML}{92400E}
\definecolor{medslate}{HTML}{334155}
\definecolor{mutedgray}{HTML}{475569}

% Κεφαλίδα και υποσέλιδο
\setlength{\headheight}{14pt}
\pagestyle{fancy}
\fancyhf{}
\renewcommand{\headrulewidth}{0pt}
\renewcommand{\footrulewidth}{0.4pt}
\renewcommand{\footrule}{\hbox to\headwidth{\color{bordergray}\leaders\hrule height \footrulewidth\hfill}}
\fancyfoot[L]{\fontsize{7.5pt}{9pt}\selectfont\color{mutedgray}Συμμόρφωση Επιχειρησιακής Διαδικασίας Ακτινοδιαγνωστικού Τμήματος Νοσοκομείου με την Οδηγία NIS2}
\fancyfoot[R]{\fontsize{7.5pt}{9pt}\selectfont\color{mutedgray}Σελίδα \thepage\ από \pageref{LastPage}}

% Ρυθμίσεις παραγράφου και αναδίπλωσης λέξεων
\setlength{\parindent}{0pt}
\setlength{\parskip}{0pt}
\emergencystretch=2em
\hyphenpenalty=50

\begin{document}

% ==================== ΕΝΟΤΗΤΑ ΤΙΤΛΟΥ ====================
{\fontsize{7.5pt}{9pt}\selectfont\color{mutedgray}\textbf{ΑΡΙΣΤΟΤΕΛΕΙΟ ΠΑΝΕΠΙΣΤΗΜΙΟ ΘΕΣΣΑΛΟΝΙΚΗΣ \textbar\ GOOGLE.ORG CYBERSECURITY SEMINARS}}\\[2pt]
{\fontsize{13pt}{16pt}\selectfont\color{primaryblue}\textbf{Ενιαίος Πίνακας Καταγραφής Περιουσιακών Στοιχείων, Κινδύνων \& Συμμόρφωσης NIS2}}\\[2pt]
{\fontsize{8.5pt}{11pt}\selectfont\color{medslate}Ζωίδης Βασίλειος — 23 Νοεμβρίου 2026}\\[1pt]
{\color{accentblue}\hrule height 1.1pt}
\vspace{3pt}

% ==================== ΕΝΙΑΙΟΣ ΠΙΝΑΚΑΣ ====================
{\fontsize{6.8pt}{8.5pt}\selectfont
\setlength{\tabcolsep}{2.3pt}
\renewcommand{\arraystretch}{1.28}
\setlength{\extrarowheight}{1.2pt}
\arrayrulecolor{bordergray}

\begin{longtable}{|
  >{\centering\arraybackslash}p{6.0mm}|
  >{\RaggedRight\arraybackslash}p{20.0mm}|
  >{\RaggedRight\arraybackslash}p{14.5mm}|
  >{\RaggedRight\arraybackslash}p{21.0mm}|
  >{\centering\arraybackslash}p{9.5mm}|
  >{\RaggedRight\arraybackslash}p{14.0mm}|
  >{\RaggedRight\arraybackslash}p{24.0mm}|
  >{\RaggedRight\arraybackslash}p{17.5mm}|
  >{\RaggedRight\arraybackslash}p{23.0mm}|
  >{\RaggedRight\arraybackslash}p{13.0mm}|
  >{\RaggedRight\arraybackslash}p{15.5mm}|
  >{\RaggedRight\arraybackslash}p{18.0mm}|
  >{\RaggedRight\arraybackslash}p{11.5mm}|
  >{\RaggedRight\arraybackslash}p{45.5mm}|}

\hline
\rowcolor{primaryblue}
\textcolor{white}{\textbf{Α/Α}} &
\textcolor{white}{\textbf{Υπηρεσία / Σύστημα}} &
\textcolor{white}{\textbf{Κατηγορία Υπηρεσίας}} &
\textcolor{white}{\textbf{Περιγραφή}} &
\textcolor{white}{\textbf{Κρίσιμη}} &
\textcolor{white}{\textbf{Owner}} &
\textcolor{white}{\textbf{Πιθανή Ευπάθεια / Απειλή}} &
\textcolor{white}{\textbf{Τεχνολογία / Πλατφόρμα}} &
\textcolor{white}{\textbf{Αξιολόγηση Κρίσιμου Κινδύνου}} &
\textcolor{white}{\textbf{Αξιολόγηση Ασφάλειας}} &
\textcolor{white}{\textbf{Πολιτική Αναφορών}} &
\textcolor{white}{\textbf{Δικτυακές Πληροφορίες}} &
\textcolor{white}{\textbf{Συσχέτιση NIS2}} &
\textcolor{white}{\textbf{Προτεινόμενη Πολιτική Ασφαλείας \& Τεχνικά Μέτρα}} \\
\hline
\endfirsthead

\hline
\rowcolor{primaryblue}
\textcolor{white}{\textbf{Α/Α}} &
\textcolor{white}{\textbf{Υπηρεσία / Σύστημα}} &
\textcolor{white}{\textbf{Κατηγορία Υπηρεσίας}} &
\textcolor{white}{\textbf{Περιγραφή}} &
\textcolor{white}{\textbf{Κρίσιμη (Ναι/Όχι)}} &
\textcolor{white}{\textbf{Owner}} &
\textcolor{white}{\textbf{Πιθανή Ευπάθεια / Απειλή}} &
\textcolor{white}{\textbf{Τεχνολογία / Πλατφόρμα}} &
\textcolor{white}{\textbf{Αξιολόγηση Κρίσιμου Κινδύνου}} &
\textcolor{white}{\textbf{Αξιολόγηση Ασφάλειας}} &
\textcolor{white}{\textbf{Πολιτική Αναφορών}} &
\textcolor{white}{\textbf{Δικτυακές Πληροφορίες}} &
\textcolor{white}{\textbf{Συσχέτιση NIS2}} &
\textcolor{white}{\textbf{Προτεινόμενη Πολιτική Ασφαλείας \& Τεχνικά Μέτρα}} \\
\hline
\endhead

\hline
\multicolumn{14}{r}{\fontsize{6.8pt}{8.5pt}\selectfont\color{mutedgray}\textit{Συνέχεια στην επόμενη σελίδα...}} \\
\endfoot

\hline
\endlastfoot

% --- 1 ---
1 &
\textbf{Firewall} (Περιμετρικό Τείχος Προστασίας) &
Υλισμικοτεχνικό / Δίκτυο &
Ενιαία προστασία περιμέτρου, διαχείριση Webservices, διασύνδεσης HIS \& Site-to-Site VPN &
\textbf{Ναι} &
ΥΑΣΠΕ &
Outdated firmware: \textbf{CVE-2023-27997} (Heap Buffer Overflow / Pre-auth RCE σε SSL-VPN), DDoS, παράκαμψη κανόνων &
Fortinet Next-Gen Firewall, \textbf{Firmware v.2.5 (2023)} &
\textcolor{critred}{\textbf{Κρίσιμος:}} Παραβίαση περιμέτρου, μη εξουσιοδοτημένη είσοδος στο εσωτερικό δίκτυο, διακοπή συνδέσεων (ΗΔΙΚΑ, Backup). &
Εξαμηνιαία \& Pen-Test &
Εσωτ.: <1ω.\newline NIS2: 24ω / 72ω / 1μ. &
Εσωτ.\ IP:\newline 10.10.20.254\newline Εξωτ.\ IP:\newline 173.82.147.219\newline 10.10.20.0/24 &
Άρθρο 21(2)(α, β, ε, η, ι), Άρθρο 23 &
\textbf{Πολιτική Συστήματος:} Άμεση αναβάθμιση firmware, απενεργοποίηση SSL-VPN εφόσον δεν απαιτείται, IPS/IDS, Geo-IP filtering, κανόνες deny-by-default, MFA στη διαχείριση, φιλτράρισμα κίνησης από HIS. \textbf{Αντίδραση:} Απομόνωση παραβιασμένων tunnels, ενεργοποίηση failover, ανάλυση logs σε SIEM. \\
\hline

% --- 2 ---
\rowcolor{rowalt}
2 &
\textbf{Δικτυακός Μεταγωγέας} (Network Switch) &
Υλισμικοτεχνικό / Δίκτυο &
Εσωτερική δικτυακή διασύνδεση servers (RIS/PACS), MRI και τερματικών Η/Υ του τμήματος &
\textbf{Ναι} &
Δ/νση Πληροφορικής &
\textbf{CVE-2023-20198} (Web UI Unauthenticated Privilege Escalation), ARP Spoofing, VLAN Hopping, MAC Flooding &
Cisco Catalyst Managed L2/L3 Switch (IOS XE) &
\textcolor{highamber}{\textbf{Υψηλός:}} Πλευρική κίνηση επιτιθέμενου από τερματικά προς MRI/Servers, υποκλοπή κυκλοφορίας DICOM. &
Ετήσια &
Εσωτ.: <4ω.\newline NIS2: 24ω / 72ω (σε διακοπή). &
IP:\newline 10.10.20.2\newline Υποδίκτυο:\newline 10.10.20.0/24 &
Άρθρο 21(2)(α, ε, θ), Άρθρο 23 &
\textbf{Πολιτική Συστήματος:} Αναβάθμιση IOS XE, απενεργοποίηση HTTP/HTTPS Web UI (διαχείριση μόνο μέσω SSHv2), προτεινόμενη κατάτμηση του 10.10.20.0/24 σε ξεχωριστά VLANs (MRI, Servers, PCs), Port Security (802.1X), κλείσιμο ανενεργών θυρών. \textbf{Αντίδραση:} Αυτόματο shutdown παραβιασμένης θύρας, απομόνωση VLAN. \\
\hline

% --- 3 ---
3 &
\textbf{Εξυπηρετητής RIS} (RIS Server) &
Υλισμικοτεχνικό / Server &
Φυσικός ή εικονικός διακομιστής φιλοξενίας της εφαρμογής και της βάσης δεδομένων RIS &
\textbf{Ναι} &
Δ/νση Πληροφορικής &
\textbf{CVE-2024-38063} (Windows TCP/IP IPv6 RCE), Ransomware, αστοχία υλικού &
Rack Server, MS Windows Server 2019 Standard &
\textcolor{critred}{\textbf{Κρίσιμος:}} Πλήρης διακοπή εγγραφών ασθενών, παραπεμπτικών και έκδοσης γνωματεύσεων· κρυπτογράφηση δεδομένων. &
Εξαμηνιαία &
Εσωτ.: <1ω.\newline NIS2: 24ω / 72ω / 1μ. &
IP:\newline 10.10.20.10\newline Υποδίκτυο:\newline 10.10.20.0/24 &
Άρθρο 21(2)(α, β, γ, ε), Άρθρο 23 &
\textbf{Πολιτική Συστήματος:} Άμεση εφαρμογή patches (και απενεργοποίηση IPv6 αν δεν χρησιμοποιείται), EDR, σκληραγώγηση OS, RAID, περιορισμός RDP μόνο από διαχειριστικό σταθμό, αμετάβλητα (immutable) backups. \textbf{Αντίδραση:} Απομόνωση host μέσω EDR, ενεργοποίηση Σχεδίου Ανάκαμψης (DRP), επαναφορά από καθαρό backup. \\
\hline

% --- 4 ---
\rowcolor{rowalt}
4 &
\textbf{Εξυπηρετητής PACS} (PACS Server \& Storage) &
Υλισμικοτεχνικό / Server &
Κεντρικός διακομιστής αποθήκευσης, αρχειοθέτησης και διανομής ιατρικών εικόνων DICOM &
\textbf{Ναι} &
Δ/νση Πληροφορικής &
\textbf{CVE-2024-6387} (OpenSSH regreSSHion RCE), Ransomware, βλάβη συστοιχίας δίσκων &
Storage Server, Red Hat Enterprise Linux 9 &
\textcolor{critred}{\textbf{Κρίσιμος:}} Απώλεια ή αδυναμία ανάκτησης διαγνωστικών εικόνων, καθυστέρηση κρίσιμων ιατρικών διαγνώσεων. &
Εξαμηνιαία &
Εσωτ.: <1ω.\newline NIS2: 24ω / 72ω / 1μ. &
IP:\newline 10.10.20.11\newline Υποδίκτυο:\newline 10.10.20.0/24 &
Άρθρο 21(2)(β, γ, ε, η), Άρθρο 23 &
\textbf{Πολιτική Συστήματος:} Ενημέρωση OpenSSH \& πυρήνα Linux, κρυπτογράφηση δίσκων (AES-256), αυτοματοποιημένα κρυπτογραφημένα immutable backups στο Remote Backup μέσω IPsec VPN, παρακολούθηση υγείας δίσκων. \textbf{Αντίδραση:} Άμεση αποσύνδεση από το δίκτυο σε υποψία ransomware, επαναφορά από Remote Backup. \\
\hline

% --- 5 ---
5 &
\textbf{Μαγνητικός Τομογράφος} (MRI Scanner) &
Υλισμικοτεχνικό / Ιατρικός Εξοπλισμός &
Εκτέλεση μαγνητικών τομογραφιών και παραγωγή πρωτογενών δεδομένων απεικόνισης (Βήμα 3) &
\textbf{Ναι} &
Προϊστάμενος Ακτινοδιαγν. &
Outdated firmware: \textbf{CVE-2021-26248} \& \textbf{CVE-2021-42744} (ελλιπής έλεγχος πρόσβασης \& έκθεση δεδομένων, CISA ICSMA-21-313-01), προεπιλεγμένοι κωδικοί, παραβίαση μέσω Vendor VPN &
MRI Scanner, \textbf{Firmware v.0.15 (2021)} &
\textcolor{critred}{\textbf{Κρίσιμος (Ασφάλεια Ασθενούς):}} Διακοπή λειτουργίας, αλλοίωση παραμέτρων σάρωσης (κυβερνοφυσικός κίνδυνος), είσοδος malware από προμηθευτή. &
Εξαμηνιαία \& Έλεγχος Προμηθευτή &
Άμεσα ΥΑΣΠΕ \& Κατασκευαστή.\newline NIS2: 24ω / 72ω / 1μ. &
IP:\newline 10.10.20.20\newline Υποδίκτυο:\newline 10.10.20.0/24 &
Άρθρο 21(2)(α, δ, ε, στ, θ), Άρθρο 23 &
\textbf{Πολιτική Συστήματος:} Αίτημα στον προμηθευτή για πιστοποιημένη αναβάθμιση firmware· εάν δεν υποστηρίζεται, αυστηρή απομόνωση σε Medical VLAN με virtual patching στο Firewall, κατάργηση default accounts, πρόσβαση Vendor VPN μόνο κατόπιν έγκρισης (Just-in-Time) με καταγραφή συνεδρίας. \textbf{Αντίδραση:} Άμεση δικτυακή απομόνωση, τεχνική επιβεβαίωση πριν τη χρήση σε ασθενείς. \\
\hline

% --- 6 ---
\rowcolor{rowalt}
6 &
\textbf{Σταθμός Εργασίας MRI} (MRI Workstation) &
Υλισμικοτεχνικό / Ιατρικό Τερματικό &
Ρύθμιση σάρωσης, προεπεξεργασία εικόνων DICOM και αποστολή στο PACS (Βήματα 3, 5) &
\textbf{Ναι} &
Προϊστάμενος Ακτινοδιαγν. &
\textbf{CVE-2023-21768} (Windows AFD WinSock EoP), μόλυνση μέσω USB, σφάλματα επεξεργασίας κακόβουλων DICOM &
Medical Workstation, MS Windows 10 Enterprise LTSC \& DICOM Acquisition SW &
\textcolor{highamber}{\textbf{Υψηλός:}} Διακοπή ροής εξετάσεων, αλλοίωση ή διαρροή ακατέργαστων ιατρικών απεικονίσεων πριν την αρχειοθέτηση. &
Ετήσια &
Εσωτ.: <4ω.\newline NIS2: 24ω / 72ω (σε σημαντικό συμβάν). &
IP:\newline 10.10.20.21\newline Υποδίκτυο:\newline 10.10.20.0/24 &
Άρθρο 21(2)(δ, ε, θ), Άρθρο 23 &
\textbf{Πολιτική Συστήματος:} Εγκεκριμένα security patches σε συνεννόηση με προμηθευτή, αποκλεισμός USB, λίστα επιτρεπόμενων εφαρμογών (allowlisting), συμβατό EDR, ατομικοί λογαριασμοί τεχνολόγων. \textbf{Αντίδραση:} Απομόνωση σταθμού, επανεγκατάσταση από πιστοποιημένο είδωλο συστήματος (golden image). \\
\hline

% --- 7 ---
7 &
\textbf{Η/Υ Ακτινοδιαγνώστη} (Radiologist PC) &
Υλισμικοτεχνικό / Τερματικό Η/Υ &
Διαγνωστική προβολή DICOM (Viewer) και έκδοση/υπογραφή γνωματεύσεων (Βήματα 7--9) &
\textbf{Ναι} &
Διευθυντής Ιατρικής Υπηρεσίας &
\textbf{CVE-2024-30051} (Windows DWM Core Library EoP), Phishing, υποκλοπή κωδικών ή token ψηφιακής υπογραφής &
MS Windows 10 Professional, Ιατρικές Οθόνες, DICOM Viewer &
\textcolor{highamber}{\textbf{Υψηλός:}} Πλαστογράφηση ιατρικών γνωματεύσεων, μη εξουσιοδοτημένη πρόσβαση σε όλους τους φακέλους ασθενών. &
Ετήσια &
Εσωτ.: <4ω.\newline NIS2: 24ω / 72ω (σε διαρροή/ransomware). &
IP:\newline 10.10.20.30\newline Υποδίκτυο:\newline 10.10.20.0/24 &
Άρθρο 21(2)(ε, ζ, η, θ, ι), Άρθρο 23 &
\textbf{Πολιτική Συστήματος:} Τακτικό patching (Windows \& DICOM Viewer), κρυπτογράφηση δίσκου (BitLocker), EDR, αφαίρεση δικαιωμάτων τοπικού διαχειριστή, αυτόματο κλείδωμα οθόνης (5 λεπτά), κρυπτογραφικό token για ψηφιακή υπογραφή. \textbf{Αντίδραση:} Ανάκληση συνεδριών/πιστοποιητικού, απομόνωση τερματικού. \\
\hline

% --- 8 ---
\rowcolor{rowalt}
8 &
\textbf{Η/Υ Εσωτερικού Ιατρού} (Internal Physician PC) &
Υλισμικοτεχνικό / Τερματικό Η/Υ &
Πρόσβαση σε γνωματεύσεις και DICOM ασθενών που εξετάζονται εντός νοσοκομείου (Βήμα 12) &
Όχι &
Διευθυντής Ιατρικής Υπηρεσίας &
\textbf{CVE-2024-38063} (Windows TCP/IP IPv6 RCE), Phishing, κοινόχρηστοι λογαριασμοί κλινικής &
MS Windows 10 Professional &
\textcolor{medslate}{\textbf{Μέτριος--Υψηλός:}} Σημείο εισόδου malware στο τμήμα, τοπική διαρροή ιατρικών δεδομένων. &
Ετήσια &
Εσωτ.: <24ω (ΕΑΚ μόνο σε σημαντικό συμβάν). &
IP:\newline 10.10.20.31\newline Υποδίκτυο:\newline 10.10.20.0/24 &
Άρθρο 21(2)(ε, ζ, θ, ι) &
\textbf{Πολιτική Συστήματος:} Τακτικό patching, web/email filtering, BitLocker, αυστηρά ατομικοί λογαριασμοί με MFA, δικαιώματα μόνο ανάγνωσης (read-only) στο RIS/PACS. \textbf{Αντίδραση:} Αποσύνδεση τερματικού, σάρωση EDR, επαναφορά κωδικών. \\
\hline

% --- 9 ---
9 &
\textbf{Η/Υ Γραμματείας \& Εγγραφέας CD} &
Υλισμικοτεχνικό / Τερματικό Η/Υ &
Εγγραφή ασθενών, ανάκτηση e-prescription, εγγραφή CD με DICOM \& παράδοση αποτελεσμάτων (Βήματα 1--3, 10--11) &
\textbf{Ναι} &
Δ/νση Διοικητικών Υπηρεσιών &
\textbf{CVE-2021-34527} (Windows Print Spooler RCE -- PrintNightmare), Phishing, κακόβουλα εξωτερικά μέσα, παραμονή προσωρινών DICOM στον δίσκο &
MS Windows 10 Professional, Optical CD/DVD Burner &
\textcolor{highamber}{\textbf{Υψηλός:}} Άμεση επαφή με κοινό \& Webservices· διαρροή προσωπικών/ιατρικών δεδομένων, αρχική μόλυνση δικτύου. &
Ετήσια &
Εσωτ.: <4ω.\newline NIS2: 24ω / 72ω (σε διαρροή ή διακοπή). &
IP:\newline 10.10.20.40\newline Υποδίκτυο:\newline 10.10.20.0/24 &
Άρθρο 21(2)(ε, ζ, η, θ), Άρθρο 23 &
\textbf{Πολιτική Συστήματος:} Πλήρες patching, περιορισμός Print Spooler, αποκλεισμός USB, αυτόματη ασφαλής διαγραφή προσωρινών αρχείων DICOM μετά την εγγραφή CD, BitLocker, φίλτρο ιδιωτικότητας οθόνης. \textbf{Αντίδραση:} Απομόνωση Η/Υ, έλεγχος εξαγόμενων αρχείων (DLP). \\
\hline

% --- 10 ---
\rowcolor{rowalt}
10 &
\textbf{Σταθμός Διαχείρισης IT} (Admin Workstation) &
Υλισμικοτεχνικό / Τερματικό IT &
Κεντρική διαχείριση Firewall, Switch, Servers (RIS/PACS), ενημερώσεων και αντιγράφων ασφαλείας &
\textbf{Ναι} &
Δ/νση Πληροφορικής &
\textbf{CVE-2024-30051} (Windows DWM EoP), Pass-the-Hash, υποκλοπή προνομιούχων διαπιστευτηρίων &
MS Windows 10 Professional (Αφιερωμένος Σταθμός Διαχείρισης) &
\textcolor{critred}{\textbf{Κρίσιμος:}} Πλήρης κατάληψη υποδομής, απενεργοποίηση backups και μαζική εγκατάσταση ransomware. &
Εξαμηνιαία &
Εσωτ.: <1ω.\newline NIS2: 24ω / 72ω / 1μ. &
IP:\newline 10.10.20.5\newline Υποδίκτυο:\newline 10.10.20.0/24 &
Άρθρο 21(2)(α, β, ε, θ, ι), Άρθρο 23 &
\textbf{Πολιτική Συστήματος:} Χρήση αποκλειστικά για διαχείριση (χωρίς internet/email), σκληραγώγηση OS, BitLocker, υποχρεωτικό MFA με hardware keys, πρόσβαση διαχείρισης servers μόνο από αυτή τη διεύθυνση IP. \textbf{Αντίδραση:} Άμεση απομόνωση, ανάκληση όλων των admin credentials, εγκληματολογικός έλεγχος. \\
\hline

% --- 11 ---
11 &
\textbf{Δικτυακός Εκτυπωτής Γνωματεύσεων} &
Υλισμικοτεχνικό / Περιφερειακό &
Εκτύπωση υπογεγραμμένων ιατρικών γνωματεύσεων προς παράδοση στον ασθενή (Βήμα 11) &
Όχι &
Δ/νση Διοικητικών Υπηρεσιών &
Προεπιλεγμένοι κωδικοί Web UI, μη κρυπτογραφημένες ουρές εκτύπωσης, παραμονή εκτυπωμένων γνωματεύσεων στον δίσκο εξόδου &
Network Laser Printer (Embedded Firmware) &
\textcolor{medslate}{\textbf{Μέτριος:}} Μη εξουσιοδοτημένη ανάγνωση έντυπων διαγνώσεων, αξιοποίηση εκτυπωτή ως σημείου εσωτερικής διείσδυσης. &
Ετήσια &
Εσωτ.: <24ω (ΕΑΚ μόνο σε σημαντικό συμβάν). &
IP:\newline 10.10.20.41\newline Υποδίκτυο:\newline 10.10.20.0/24 &
Άρθρο 21(2)(ε, θ) &
\textbf{Πολιτική Συστήματος:} Ενημέρωση firmware εκτυπωτή, αλλαγή default passwords, απενεργοποίηση αχρησιμοποίητων πρωτοκόλλων (Telnet/FTP), ασφαλής εκτύπωση με PIN/κάρτα (Follow-Me Print), τοποθέτηση σε ελεγχόμενο χώρο γραμματείας. \textbf{Αντίδραση:} Αποσύνδεση από το δίκτυο, επαναφορά ρυθμίσεων. \\
\hline

% --- 12 ---
\rowcolor{rowalt}
12 &
\textbf{Σύστημα UPS} (Αδιάλειπτη Τροφοδοσία) &
Υλισμικοτεχνικό / Κρίσιμη Υποδομή &
Αδιάλειπτη παροχή ενέργειας σε RIS/PACS Servers, Firewall, Switch \& MRI Workstation &
\textbf{Ναι} &
Τεχνική Υπηρεσία Νοσοκομείου &
Διακοπή ηλεκτροδότησης, εξάντληση συσσωρευτών, εκμετάλλευση μη ασφαλούς κάρτας διαχείρισης (SNMP v1/v2c) &
Online Double-Conversion UPS, SNMP Management Card &
\textcolor{highamber}{\textbf{Υψηλός:}} Αιφνίδια διακοπή λειτουργίας servers \& αλλοίωση βάσεων δεδομένων, παύση εξυπηρέτησης ασθενών. &
Ετήσια \& Δοκιμή Φορτίου &
Εσωτ.: <1ω.\newline NIS2: 24ω / 72ω (σε παρατεταμένη διακοπή). &
IP:\newline 10.10.20.6\newline Υποδίκτυο:\newline 10.10.20.0/24 &
Άρθρο 21(2)(γ, ε), Άρθρο 23 &
\textbf{Πολιτική Συστήματος:} Μετάβαση σε SNMPv3 με κρυπτογράφηση, απομόνωση κάρτας διαχείρισης σε Management VLAN, σύνδεση με εφεδρική ηλεκτρογεννήτρια νοσοκομείου, προληπτική συντήρηση μπαταριών, ειδοποιήσεις αυτόματου τερματισμού servers. \textbf{Αντίδραση:} Ενεργοποίηση γεννήτριας \& σχεδίου επιχειρησιακής συνέχειας. \\
\hline

% --- 13 ---
13 &
\textbf{Πληροφοριακό Σύστημα RIS \& Βάση Δεδομένων} &
Λογισμικοτεχνικό / Πληρ.\ Σύστημα &
Διαχείριση ροής ασθενών, παραπεμπτικών e-prescription, ραντεβού και γνωματεύσεων (Βήματα 2--3, 7, 9--10, 12) &
\textbf{Ναι} &
Δ/νση Πληροφορικής &
SQL Injection, Broken Access Control (OWASP), αδύναμη αυθεντικοποίηση, κατάχρηση API Webservices (ΗΔΙΚΑ/ΕΟΠΥΥ) ή διασύνδεσης HIS &
RIS Application, MS SQL Server, HL7 / REST Webservices &
\textcolor{critred}{\textbf{Κρίσιμος:}} Μαζική διαρροή ιατρικών ιστορικών, αλλοίωση γνωματεύσεων, κατάχρηση πρόσβασης προς εθνικά συστήματα υγείας. &
Εξαμηνιαία (Vulnerability Scan) &
Εσωτ.: <1ω.\newline NIS2: 24ω / 72ω / 1μ. &
Φιλοξενία:\newline 10.10.20.10\newline (10.10.20.0/24) &
Άρθρο 21(2)(α, β, γ, ε, η, θ, ι), Άρθρο 23 &
\textbf{Πολιτική Συστήματος:} Κρυπτογράφηση βάσης (TDE) \& επικοινωνιών (TLS 1.3), αυστηρό RBAC \& MFA, παραμετροποιημένα ερωτήματα (αποτροπή SQLi), WAF, αμοιβαία πιστοποιητικά (mTLS) για Webservices, πλήρες audit logging σε SIEM. \textbf{Αντίδραση:} Αναστολή παραβιασμένων λογαριασμών, ανάλυση logs, επαναφορά βάσης. \\
\hline

% --- 14 ---
\rowcolor{rowalt}
14 &
\textbf{Σύστημα PACS \& Βάση Δεδομένων} &
Λογισμικοτεχνικό / Πληρ.\ Σύστημα &
Αποθήκευση, ευρετηρίαση, ανάκτηση και απομακρυσμένο backup ιατρικών εικόνων DICOM (Βήματα 4--8) &
\textbf{Ναι} &
Δ/νση Πληροφορικής &
Μη κρυπτογραφημένο πρωτόκολλο DICOM (cleartext), AE Title spoofing, SQL Injection, αποτυχία αντιγράφων ασφαλείας &
PACS DICOM Server, PostgreSQL, IPsec Backup Agent &
\textcolor{critred}{\textbf{Κρίσιμος:}} Υποκλοπή ή σκόπιμη αλλοίωση διαγνωστικών εικόνων (κίνδυνος εσφαλμένης διάγνωσης), διακοπή διαγνωστικής λειτουργίας. &
Εξαμηνιαία &
Εσωτ.: <1ω.\newline NIS2: 24ω / 72ω / 1μ. &
Φιλοξενία:\newline 10.10.20.11\newline (10.10.20.0/24) &
Άρθρο 21(2)(β, γ, ε, η, θ), Άρθρο 23 &
\textbf{Πολιτική Συστήματος:} Ενεργοποίηση DICOM over TLS, λίστα επιτρεπόμενων AE Titles \& IP whitelisting, κρυπτογράφηση δεδομένων σε ηρεμία (AES-256), τριμηνιαίες δοκιμές επαναφοράς από Remote Backup. \textbf{Αντίδραση:} Αποκλεισμός μη εξουσιοδοτημένων DICOM nodes, έλεγχος ακεραιότητας εικόνων (checksums), επαναφορά backup. \\
\hline

% --- 15 ---
15 &
\textbf{Υπηρεσία Τερματισμού VPN \& Webservices} (Εφοδιαστική Αλυσίδα) &
Λογισμικοτεχνικό / Διασύνδεση &
Εσωτερικά τερματικά σημεία επικοινωνίας με Προμηθευτή MRI, Remote Backup, ΗΔΙΚΑ/ΕΟΠΥΥ, Εξωτ.\ Ιατρό \& HIS &
\textbf{Ναι} &
ΥΑΣΠΕ &
Επίθεση εφοδιαστικής αλυσίδας (Supply Chain), διαρροή κλειδιών VPN/API, πλευρική κίνηση από προμηθευτή ή άλλα τμήματα (HIS) &
Site-to-Site IPsec VPN, TLS 1.3 Webservices Gateway &
\textcolor{critred}{\textbf{Κρίσιμος:}} Είσοδος επιτιθέμενου μέσω αξιόπιστου τρίτου (προμηθευτής MRI/HIS), διακοπή εξωτερικής επικοινωνίας \& backup. &
Εξαμηνιαία \& Έλεγχος SLA &
Εσωτ.: <1ω.\newline NIS2: 24ω / 72ω / 1μ. &
Τερματισμός:\newline 10.10.20.254\newline / 173.82.147.219 &
Άρθρο 21(2)(α, β, δ, η, ι), Άρθρο 23 &
\textbf{Πολιτική Συστήματος:} Ρήτρες κυβερνοασφάλειας (SLA) με προμηθευτές MRI \& Backup, IPsec IKEv2/AES-256 με περιοδική εναλλαγή κλειδιών, μόνιμα απενεργοποιημένο Vendor VPN (ενεργοποίηση μόνο κατόπιν αιτήματος), περιορισμός πρόσβασης μόνο στην IP του MRI, απομόνωση κίνησης HIS. \textbf{Αντίδραση:} Άμεσος τερματισμός VPN tunnel / API tokens, ειδοποίηση προμηθευτή. \\
\hline

% --- 16 ---
\rowcolor{rowalt}
16 &
\textbf{Δεδομένα Ασθενών \& Παραπεμπτικών} (e-Prescription) &
Πληροφοριακά Δεδομένα &
Προσωπικά, ασφαλιστικά και κλινικά δεδομένα εξωτερικών ασθενών (Βήματα 1--3) &
\textbf{Ναι} &
Διοίκηση Νοσοκομείου &
Διαρροή ειδικών κατηγοριών δεδομένων υγείας, αλλοίωση δημογραφικών, υποκλοπή κατά τη μεταφορά από ΗΔΙΚΑ/ΕΟΠΥΥ &
Βάση RIS (SQL), XML/JSON (Webservices) &
\textcolor{critred}{\textbf{Κρίσιμος:}} Παραβίαση ιατρικού απορρήτου, κυρώσεις (NIS2 \& GDPR), κίνδυνος ταυτοποίησης λανθασμένου ασθενούς. &
Ετήσια &
Εσωτ.: <1ω.\newline NIS2: 24ω / 72ω / 1μ (συν ΑΠΔΠΧ). &
Δ/Ε\newline (Φιλοξενία:\newline 10.10.20.10) &
Άρθρο 21(2)(α, β, η, θ), Άρθρο 23 &
\textbf{Πολιτική Συστήματος:} Ταξινόμηση ως «Αυστηρά Εμπιστευτικά», ελαχιστοποίηση δεδομένων, κρυπτογράφηση at-rest \& in-transit (TLS 1.3), πρόληψη διαρροών (DLP), πολιτική διατήρησης \& ασφαλούς διαγραφής. \textbf{Αντίδραση:} Ενεργοποίηση διαδικασίας παραβίασης δεδομένων, περιορισμός διαρροής, κοινοποίηση στις αρχές. \\
\hline

% --- 17 ---
17 &
\textbf{Ιατρικές Απεικονίσεις DICOM} (Ενεργές \& Backups) &
Πληροφοριακά Δεδομένα &
Διαγνωστικές εικόνες MRI (με metadata ασθενούς) σε PACS και Remote Backup (Βήματα 4--8) &
\textbf{Ναι} &
Διευθυντής Ακτινοδιαγν. &
Υποκλοπή, κρυπτογράφηση ransomware, αλλοίωση DICOM headers, απώλεια λόγω αστοχίας backup &
Πρότυπο DICOM (.dcm), Κρυπτογραφημένο IPsec Backup &
\textcolor{critred}{\textbf{Κρίσιμος:}} Απώλεια μοναδικών διαγνωστικών εξετάσεων (ανάγκη επανάληψης σάρωσης), παραβίαση απορρήτου. &
Ετήσια &
Εσωτ.: <1ω.\newline NIS2: 24ω / 72ω / 1μ. &
Δ/Ε\newline (Φιλοξενία:\newline 10.10.20.11,\newline 10.10.20.21) &
Άρθρο 21(2)(γ, η, θ), Άρθρο 23 &
\textbf{Πολιτική Συστήματος:} Κρυπτογράφηση αποθήκευσης \& μεταφοράς (DICOM-TLS, IPsec), κανόνας αντιγράφων 3-2-1 με αμετάβλητα backups, ψηφιακά checksums ακεραιότητας, ανωνυμοποίηση σε τυχόν εκπαιδευτική χρήση. \textbf{Αντίδραση:} Επαναφορά από αμετάβλητο backup, διερεύνηση εύρους παραβίασης. \\
\hline

% --- 18 ---
\rowcolor{rowalt}
18 &
\textbf{Υπογεγραμμένες Γνωματεύσεις \& Φυσικά CD} &
Πληροφοριακά Δεδομένα &
Τελικές ιατρικές διαγνώσεις και εξαγόμενα CD προς ασθενείς ή εξωτερικούς ιατρούς (Βήματα 9--12) &
\textbf{Ναι} &
Διευθυντής Ακτινοδιαγν. &
Απώλεια/κλοπή μη κρυπτογραφημένου CD, πλαστογράφηση γνωμάτευσης, παράδοση σε λανθασμένο πρόσωπο &
Ψηφιακά Υπογεγραμμένα PDF (PKI), CD-R (DICOMDIR) &
\textcolor{highamber}{\textbf{Υψηλός:}} Παραβίαση ιατρικού απορρήτου εκτός περιμέτρου νοσοκομείου, νομική ευθύνη από μη εξουσιοδοτημένη διάθεση. &
Ετήσια &
Εσωτ.: Άμεσα.\newline NIS2: 24ω / 72ω (σε συστημική διαρροή). &
Δ/Ε\newline (Εξαγωγή:\newline 10.10.20.10,\newline 10.10.20.40) &
Άρθρο 21(2)(η, θ, ι), Άρθρο 23 &
\textbf{Πολιτική Συστήματος:} Εγκεκριμένη ψηφιακή υπογραφή (PKI) στις γνωματεύσεις, προστασία CD με κωδικό ή μετάβαση σε ασφαλή πύλη ασθενών (Patient Portal με MFA), αυστηρή ταυτοποίηση κατά την παράδοση. \textbf{Αντίδραση:} Καταγραφή περιστατικού, ανάκληση τυχόν συνδέσμων πρόσβασης. \\
\hline

% --- 19 ---
19 &
\textbf{Αρχεία Καταγραφής (Audit Logs) \& Συγκαταθέσεις} &
Πληροφοριακά Δεδομένα &
Καταγραφές πρόσβασης/συμβάντων συστημάτων και εγκρίσεις ασθενών για εξωτερικούς ιατρούς (Βήμα 12) &
\textbf{Ναι} &
ΥΑΣΠΕ &
Διαγραφή ή αλλοίωση logs από επιτιθέμενο (anti-forensics), διάθεση δεδομένων σε εξωτερικό ιατρό χωρίς τεκμηριωμένη συγκατάθεση &
Κεντρική Βάση Logs (Syslog/SIEM), Ψηφιακά Αρχεία Συγκατάθεσης (RIS) &
\textcolor{highamber}{\textbf{Υψηλός:}} Αδυναμία ανίχνευσης και διερεύνησης κυβερνοεπιθέσεων, παραβίαση νομιμότητας διάθεσης δεδομένων (Βήμα 12). &
Εξαμηνιαία &
Εσωτ.: <1ω.\newline NIS2: 24ω / 72ω (σε παραποίηση logs). &
Δ/Ε\newline (Φιλοξενία:\newline 10.10.20.10) &
Άρθρο 21(2)(α, β, στ, θ), Άρθρο 23 &
\textbf{Πολιτική Συστήματος:} Κεντρική συλλογή logs σε αμετάβλητη αποθήκευση (write-once) με συγχρονισμό ώρας (NTP), διατήρηση κατ' ελάχιστον 12 μήνες, υποχρεωτική ηλεκτρονική επαλήθευση \& καταγραφή συγκατάθεσης ασθενούς πριν την πρόσβαση εξωτερικού ιατρού. \textbf{Αντίδραση:} Διατήρηση αποδεικτικών στοιχείων, διερεύνηση από ΥΑΣΠΕ. \\
\hline

% --- 20 ---
\rowcolor{rowalt}
20 &
\textbf{Διοικητικός Υπάλληλος / Γραμματεία} &
Ανθρώπινοι Πόροι &
Υποδοχή ασθενών, διαχείριση e-prescription, εγγραφή RIS, εγγραφή/παράδοση CD (Βήματα 1--3, 10--11) &
\textbf{Ναι} &
Δ/νση Διοικητικών Υπηρεσιών &
Κοινωνική μηχανική (Phishing / Pretexting), ανθρώπινο λάθος (λάθος φάκελος), εσωτερική απειλή, αδύναμοι κωδικοί &
Χρήστης RIS, e-Prescription Webservices, MS Windows 10 &
\textcolor{highamber}{\textbf{Υψηλός:}} Διαρροή δεδομένων ασθενών ή αρχική μόλυνση δικτύου μέσω κακόβουλου email/συνδέσμου. &
Ετήσια Εκπαίδευση &
Εσωτ.: Άμεσα (<1ω) στον ΥΑΣΠΕ. &
Δ/Ε\newline (Πρόσβαση:\newline 10.10.20.40) &
Άρθρο 20(2), Άρθρο 21(2)(ζ, θ) &
\textbf{Πολιτική Συστήματος:} Ετήσια εκπαίδευση κυβερνοϋγιεινής \& ασκήσεις phishing, αρχή ελάχιστων προνομίων (πρόσβαση μόνο σε διοικητικά πεδία), ρήτρες εμπιστευτικότητας, τυποποιημένη διαδικασία ταυτοποίησης ασθενούς. \textbf{Αντίδραση:} Άμεσο κλείδωμα λογαριασμού σε υποψία παραβίασης, στοχευμένη επανεκπαίδευση. \\
\hline

% --- 21 ---
21 &
\textbf{Τεχνολόγος Ακτινολόγος} &
Ανθρώπινοι Πόροι &
Ρύθμιση MRI Scanner, εκτέλεση εξετάσεων και επεξεργασία/αποστολή DICOM (Βήματα 3, 5) &
\textbf{Ναι} &
Προϊστάμενος Ακτινοδιαγν. &
Λανθασμένη παραμετροποίηση τομογράφου, κοινόχρηστοι κωδικοί μεταξύ βαρδιών, σύνδεση μη εγκεκριμένων USB &
Χειριστής MRI Scanner, MRI Workstation, PACS &
\textcolor{highamber}{\textbf{Υψηλός:}} Εισαγωγή malware σε κρίσιμο ιατρικό εξοπλισμό, αδυναμία απόδοσης ευθύνης σε περιστατικό. &
Ετήσια Εκπαίδευση &
Άμεσα σε ΥΑΣΠΕ \& Βιοϊατρική Τεχνολογία. &
Δ/Ε\newline (Πρόσβαση:\newline 10.10.20.20,\newline 10.10.20.21) &
Άρθρο 21(2)(δ, ζ, θ) &
\textbf{Πολιτική Συστήματος:} Ατομικοί λογαριασμοί πρόσβασης (απαγόρευση shared accounts), εκπαίδευση σε ασφάλεια ιατροτεχνολογικών συσκευών και επιτήρηση τεχνικών προμηθευτή, αυστηρή απαγόρευση προσωπικών USB. \textbf{Αντίδραση:} Αναστολή χρήσης εξοπλισμού σε ανωμαλία λειτουργίας, ενημέρωση ΥΑΣΠΕ. \\
\hline

% --- 22 ---
\rowcolor{rowalt}
22 &
\textbf{Ακτινοδιαγνώστης Ιατρός} &
Ανθρώπινοι Πόροι &
Ανάκτηση DICOM, κλινική αξιολόγηση και ψηφιακή υπογραφή γνωματεύσεων (Βήματα 7--9) &
\textbf{Ναι} &
Διευθυντής Ιατρικής Υπηρεσίας &
Υποκλοπή κωδικών ή token ψηφιακής υπογραφής, παραμονή ξεκλείδωτου σταθμού, στοχευμένο phishing &
Χρήστης RIS/PACS, DICOM Viewer, PKI e-Token &
\textcolor{highamber}{\textbf{Υψηλός:}} Παραποίηση ιατρικών γνωματεύσεων, μαζική εξαγωγή διαγνωστικών εικόνων και ευαίσθητων δεδομένων. &
Ετήσια Εκπαίδευση &
Άμεσα (<1ω) σε απώλεια token ή ύποπτη ενέργεια. &
Δ/Ε\newline (Πρόσβαση:\newline 10.10.20.30) &
Άρθρο 21(2)(ζ, η, θ, ι) &
\textbf{Πολιτική Συστήματος:} Ισχυρή αυθεντικοποίηση (MFA / Smartcard), ασφαλής φύλαξη κρυπτογραφικού token υπογραφής, εκπαίδευση κυβερνοασφάλειας, αυτόματες ειδοποιήσεις για ανώμαλες μαζικές προσβάσεις σε φακέλους. \textbf{Αντίδραση:} Άμεση ανάκληση ψηφιακού πιστοποιητικού και κλείδωμα λογαριασμού. \\
\hline

% --- 23 ---
23 &
\textbf{Εσωτερικός Ιατρός Νοσοκομείου} &
Ανθρώπινοι Πόροι &
Αίτημα και προβολή γνωματεύσεων/DICOM για ασθενείς του νοσοκομείου (Βήμα 12) &
Όχι &
Διευθυντής Ιατρικής Υπηρεσίας &
Μη εξουσιοδοτημένη πρόσβαση σε φακέλους μη θεραπευόμενων ασθενών (snooping), κοινόχρηστοι λογαριασμοί &
Χρήστης RIS/PACS \& Πληροφοριακού Συστήματος Νοσοκομείου &
\textcolor{medslate}{\textbf{Μέτριος:}} Εσωτερική παραβίαση ιατρικού απορρήτου, κατάχρηση δικαιωμάτων πρόσβασης. &
Ετήσια &
Εσωτ.: <24ω σε ΥΑΣΠΕ (και DPO). &
Δ/Ε\newline (Πρόσβαση:\newline 10.10.20.31) &
Άρθρο 21(2)(ζ, θ) &
\textbf{Πολιτική Συστήματος:} Πολιτική «Need-to-Know» (πρόσβαση μόνο σε ανατεθειμένους ασθενείς με μηχανισμό έκτακτης πρόσβασης «Break-the-Glass»), δειγματοληπτικός έλεγχος access logs, ευαισθητοποίηση σε θέματα απορρήτου. \textbf{Αντίδραση:} Διερεύνηση συμβάντος, ανάκληση δικαιωμάτων πρόσβασης. \\
\hline

% --- 24 ---
\rowcolor{rowalt}
24 &
\textbf{Διαχειριστές Συστημάτων / Δικτύων \& ΥΑΣΠΕ} &
Ανθρώπινοι Πόροι &
Διαχείριση υποδομών, ενημερώσεων, αντιγράφων ασφαλείας, περιστατικών και συμμόρφωσης NIS2 &
\textbf{Ναι} &
Διοίκηση Νοσοκομείου &
Παραβίαση προνομιούχων λογαριασμών (Admin/Root), εσφαλμένη παραμετροποίηση Firewall/Backup, έλλειψη διαχωρισμού καθηκόντων &
Σταθμός Διαχείρισης IT, Active Directory, SIEM, PAM &
\textcolor{critred}{\textbf{Κρίσιμος:}} Πλήρης έλεγχος ολόκληρης της υποδομής από επιτιθέμενο σε περίπτωση υποκλοπής admin credentials. &
Εξαμηνιαία \& Έλεγχος Δικαιωμάτων &
Άμεση ενεργοποίηση ομάδας κρίσεων.\newline NIS2: 24ω / 72ω / 1μ. &
Δ/Ε\newline (Πρόσβαση:\newline 10.10.20.5) &
Άρθρο 20, Άρθρο 21(2)(α, β, στ, ζ, θ, ι), Άρθρο 23 &
\textbf{Πολιτική Συστήματος:} Διαχείριση Προνομιούχων Λογαριασμών (PAM), υποχρεωτικό MFA με hardware keys, διπλή έγκριση (2 διαχειριστές) σε κρίσιμες αλλαγές, συνεχής κατάρτιση NIS2 \& ασκήσεις διαχείρισης περιστατικών. \textbf{Αντίδραση:} Άμεση ανάκληση προνομιούχων tokens, εκτέλεση Σχεδίου Απόκρισης (IRP). \\
\hline

\end{longtable}
}

% ==================== ΥΠΟΜΝΗΜΑ NIS2 ====================
\vspace{2pt}
{\fontsize{7.2pt}{9.2pt}\selectfont
\setlength{\tabcolsep}{5pt}
\renewcommand{\arraystretch}{1.15}
\arrayrulecolor{bordergray}
\begin{tabular}{|>{\RaggedRight\arraybackslash}p{155mm}|>{\RaggedRight\arraybackslash}p{111mm}|}
\hline
\rowcolor{rowalt}
\textbf{\color{primaryblue}Υπόμνημα Άρθρων Οδηγίας NIS2 (ΕΕ 2022/2555 / Ν.\ 5160/2024)} &
\textbf{\color{primaryblue}Συντομογραφίες Προθεσμιών Αναφοράς (Άρθρο 23) \& Δικτύου} \\
\hline
\textbf{Άρθρο 20:} Διακυβέρνηση, έγκριση μέτρων από τη Διοίκηση \& υποχρεωτική κατάρτιση.\newline
\textbf{Άρθρο 21(2):} Μέτρα διαχείρισης κινδύνων: \textbf{(α)} Ανάλυση κινδύνου \& ασφάλεια ΠΣ, \textbf{(β)} Διαχείριση περιστατικών, \textbf{(γ)} Επιχειρησιακή συνέχεια, backups \& ανάκαμψη, \textbf{(δ)} Ασφάλεια εφοδιαστικής αλυσίδας, \textbf{(ε)} Ασφάλεια απόκτησης/ανάπτυξης/συντήρησης \& διαχείριση ευπαθειών, \textbf{(στ)} Αξιολόγηση αποτελεσματικότητας μέτρων, \textbf{(ζ)} Βασικές πρακτικές κυβερνοϋγιεινής \& εκπαίδευση, \textbf{(η)} Πολιτικές κρυπτογραφίας/κρυπτογράφησης, \textbf{(θ)} Ασφάλεια ανθρώπινου δυναμικού, έλεγχος πρόσβασης \& διαχείριση assets, \textbf{(ι)} Πολυπαραγοντική αυθεντικοποίηση (MFA) \& ασφαλείς επικοινωνίες. &
\textbf{24ω:} Έγκαιρη προειδοποίηση προς Εθνική Αρχή Κυβερνοασφάλειας (ΕΑΚ) / CSIRT για σημαντικό περιστατικό.\newline
\textbf{72ω:} Επίσημη κοινοποίηση περιστατικού με αρχική αξιολόγηση σοβαρότητας \& επιπτώσεων.\newline
\textbf{1μ:} Τελική αναλυτική έκθεση αιτίων και μέτρων μετριασμού.\newline
\textbf{Εσωτ.:} Εσωτερική κλιμάκωση στον ΥΑΣΠΕ (Υπεύθυνος Ασφάλειας Συστημάτων Πληροφορικής \& Επικοινωνιών).\newline
\textbf{Δ/Ε:} Δεν εφαρμόζεται (μη δικτυακό περιουσιακό στοιχείο — αναγράφεται ενδεικτικά ο σταθμός/εξυπηρετητής πρόσβασης). \\
\hline
\end{tabular}
}

\end{document}